\section{Lecture 13 (October 13th)}
\begin{thm}
(Hund's rule) The Hund's rule determines the ground state of atoms with $Z\leq 40$. In this scenario, the coloumb electron-electron interactions supersede the spin-orbit interactions and we can use the Russel-Saunders scheme. The rules are that
\begin{itemize}
\item[(i)] We need to maximize spin
\item[(ii)] Maximise orbital angular momentum given the spin above
\item[(iii)] If the spins are less than half filled, 
\item[(iv)] In all the processes above, we need to keep in mind the Pauli-exclusion principle
\end{itemize}
\end{thm}
\vspace{2ex}
\begin{ex}
(Carbon $Z=6$ and $Z=8$) The electron configuration is given by $(1s)^2(2s)^2(2p)^2$ and maximising spin gives $S=1$ and $l=1$. We thus have
\[{}^{3}P_{0}\]
with $2s+1=3$ and $j=0$. The electron configuration is given by $(1s)^2(2s)^2(2p)^{4}$ and maximising spin gives $S=1$ and $l=1$. We thus have
\[{}^{3}P_{2}\]
\end{ex}
\vspace{2ex}
\begin{ex}
(Maganese $Z=25$) The electron configuration is given by 
\[{}^{6}S_{5/2}\]
\end{ex}
\vspace{2ex}

